\chapter{Benchmark Methodology}
\label{ch:benchmark_methodology}

\begin{itemize}
    \item Benchmark goals
    \begin{itemize}
        \item Measure total runtime effect.
        \item Separate locality benefit from sorting overhead.
        \item Identify which sorting variants are promising.
        \item Determine whether SoA and neighbor-list sorting make the sorted order useful.
    \end{itemize}

    \item Benchmark environment
    \begin{itemize}
        \item CoolMuc.
        \item Intel Xeon Platinum 8380.
        \item Single-thread execution.
        \item OpenMP restricted to one thread.
        \item VTK disabled.
        \item Progress output disabled for benchmark runs.
        \item Metadata recorded for reproducibility.
    \end{itemize}

    \item Benchmark scenarios
    \begin{itemize}
        \item Representative homogeneous grid
        \begin{itemize}
            \item Main scenario for comparing all variants.
        \end{itemize}
        \item Fixed-density increasing domain size
        \begin{itemize}
            \item Increase particle count while keeping local density fixed.
            \item Tests scaling with working-set size.
        \end{itemize}
        \item Fixed-domain increasing density
        \begin{itemize}
            \item Keep physical domain size fixed.
            \item Increase particle count by reducing spacing.
            \item Tests effect of longer neighbor lists and higher interaction pressure.
        \end{itemize}
        \item AoS versus SoA variants
        \begin{itemize}
            \item Tests whether sorted traversal order is actually exploited.
        \end{itemize}
    \end{itemize}

    \item Variants
    \begin{itemize}
        \item No-sorting baseline.
        \item Cell linear, Morton, Hilbert.
        \item Block linear, Morton, Hilbert.
        \item Particle linear, Morton, Hilbert.
        \item Additional SoA and neighbor-list sorting configurations if benchmarked.
    \end{itemize}

    \item Runtime measurement
    \begin{itemize}
        \item Repeated simulation runs.
        \item Mean runtime.
        \item Minimum and maximum runtime.
        \item Standard deviation.
        \item Speedup relative to no sorting.
    \end{itemize}

    \item Runtime decomposition
    \begin{itemize}
        \item VTune Hotspots top-down reports.
        \item Force traversal.
        \item Neighbor-list construction.
        \item Sorting and key generation.
        \item Reference rebuild / container update.
        \item Other runtime.
        \item Used to decide whether speedups come from the intended phase.
    \end{itemize}

    \item LIKWID measurements
    \begin{itemize}
        \item \texttt{L2CACHE}
        \begin{itemize}
            \item L2 miss ratio.
            \item L2 miss rate.
            \item Primary cache-locality metric.
        \end{itemize}
        \item \texttt{CYCLE\_STALLS}
        \begin{itemize}
            \item L1D miss stall rate.
            \item L2 miss stall rate.
            \item Memory-load stall rate.
            \item Shows whether cache misses affect execution.
        \end{itemize}
        \item \texttt{L3}
        \begin{itemize}
            \item L3 data volume.
            \item L3 bandwidth.
            \item Supporting shared-cache traffic metric.
        \end{itemize}
    \end{itemize}

    \item Metric caveats
    \begin{itemize}
        \item No direct reliable L3 miss ratio available.
        \item L3 is shared and can be affected by other jobs.
        \item LIKWID measurements are whole-run, not function-region measurements.
        \item Since force traversal dominates runtime, whole-run locality metrics are still informative.
        \item Region-specific LIKWID markers are future work.
    \end{itemize}
\end{itemize}
